\documentclass[a4paper]{article}     
\usepackage{amsfonts}
\usepackage{amsmath}
\usepackage{amssymb}
\usepackage{graphicx}
\usepackage{cancel}

\newcommand{\asa}{\Leftrightarrow} %als slechts als
\newcommand{\adR}{\Rightarrow} %als dan Rechts
\newcommand{\adL}{\Leftarrow}
\newcommand{\Lh}{\left(} %Links haakje
\newcommand{\Rh}{\right)}
\newcommand{\Lv}{\left[} %Links vierkanthaakje
\newcommand{\Rv}{\right]}
\newcommand{\La}{\left\{} %Links acolade
\newcommand{\Ra}{\right\}}
\newcommand{\Ras}{\right.} %Rechts acolade stelsel (omdat om stelsels te maken heb je een punt nodig
\newcommand{\pnR}{\rightarrow} %pijl naar Rechts
\newcommand{\limiet}{\lim\limits_}
\newcommand{\schrap}{\cancel}
\newcommand{\schrapb}{\bcancel} %backslash
\newcommand{\schrapk}{\xcancel} %kruis
\newcommand{\vervang}{\cancelto}
\newcommand{\intgrl}{\displaystyle\int} %integraal teken (bij het berekenen van primitieven)

\begin{document}

\noindent \large Auteur: Yoren Collier \\
\noindent \large Studierichting: Informatica\\
\noindent \large Oefeningenreeks 4A nr. 18.14\\

\medskip

\normalsize

$ I =  \intgrl x \sin \Lh x \Rh \cos \Lh x \Rh dx $

$\asa I = \intgrl\dfrac{1}{2}x \sin \Lh 2x \Rh dx = \dfrac{1}{2} \intgrl x \sin \Lh 2x \Rh d$

$\La
\begin{array}{l}
u = x dx  \\
dv = \sin \Lh 2x \Rh  
\end{array}
\Ras
\adR
\La
\begin{array}{l}
du = 1 dx  \\
v = \dfrac{-1}{2} \cos \Lh 2x \Rh  
\end{array}
\Ras$\\

Dit waarbij je van $dv$ naar $v$ gaat via substitutie met $t = 2x$.\\

$I = \dfrac{1}{2} \Lv \dfrac{-1}{2}x \cos \Lh 2x \Rh - \intgrl \dfrac{1}{2} \cos \Lh 2x \Rh \cdot 1 dx \Rv = \dfrac{1}{2} \Lv \dfrac{-1}{2}x \cos \Lh 2x \Rh + \dfrac{1}{2} \intgrl \cos \Lh 2x \Rh dx \Rv$\\

$= \dfrac{-1}{4}x \cos \Lh 2x \Rh + \dfrac{1}{8} \sin \Lh 2x \Rh +c $\\

\begin{thebibliography}{999}
%\bibitem{a} Referentie
\end{thebibliography}
\end{document} 